% 部件（src/engine/hit-parts.dtx）：同一份主文件把全部部件写上，每个分支照 iota 的规矩自己取舍。
% 学位论文专有的部件在报告里不排，也不报错；答辩页与声明博士后没有，永久通讯地址只有博士后有。
% 前置、正文、后置由部件与第一章自动切换，三个 matter 命令都不写。标题两套写法：学位论文用 \chapter，
% 报告用 \section（报告的顶级是节）。
%! files: achievements-no-year.bib
%! allow: 章节没给英文名
{\hitif[lang=en]{\bilingualswapon}{\bilingualswapoff}
\hitsetup{info={title={部件测试}[Parts test],keywords={甲,乙}[{a,b}]}}}
\begin{document}
\cover
\titlepage
\begin{abstract}中文摘要正文。\end{abstract}
\begin{abstract}[lang=en]English abstract body.\end{abstract}
\begin{symbols}
\item[$p$] 压力，Pa
\end{symbols}
\listofabbreviations
\tableofcontents
\hitif[stage=final]{\chapter{绪论}[Introduction]}{\section{绪论}}
正文。
\begin{conclusions}结论正文。\end{conclusions}
\begin{achievements}[achievements-no-year]
  \begin{pubgroup}{参与的科研项目及获奖情况}
    \item{pub:proj}
  \end{pubgroup}
\end{achievements}
\defense
\declarations
\begin{acknowledgements}致谢正文。\end{acknowledgements}
\begin{resume}简历正文。\end{resume}
\begin{permanentaddress}地址正文。\end{permanentaddress}
\end{document}
